\documentclass[11pt]{article}

\usepackage[T1]{fontenc}
\usepackage{lmodern}
\usepackage{microtype}
\usepackage{amsmath,amssymb,amsthm,mathtools}
\numberwithin{equation}{section}
\usepackage{enumitem}
\usepackage{array}
\usepackage{tabularx}
\usepackage{geometry}
\geometry{margin=1in}
\usepackage{cite}
\usepackage{hyperref}
\usepackage{url}
\usepackage{tikz}
\hypersetup{colorlinks=true,linkcolor=blue,citecolor=blue,urlcolor=blue}
\setlength{\emergencystretch}{3em}

\theoremstyle{plain}
\newtheorem{theorem}{Theorem}[section]
\newtheorem{lemma}[theorem]{Lemma}
\newtheorem{proposition}[theorem]{Proposition}
\newtheorem{corollary}[theorem]{Corollary}
\theoremstyle{definition}
\newtheorem{definition}[theorem]{Definition}
\theoremstyle{remark}
\newtheorem{remark}[theorem]{Remark}
\newtheorem{example}[theorem]{Example}
\newtheorem{openproblem}[theorem]{Open Problem}

\newcommand{\Wit}{\mathsf{W}}
\newcommand{\Cmcf}[2]{\mathcal{C}^{\mathrm{mcf}}_{#1,#2}}
\newcommand{\Cboundedtyping}[2]{\mathcal{C}^{\mathrm{mcf},\le #2}_{#1}}
\newcommand{\Calltyping}[1]{\mathcal{C}^{\mathrm{mcf},\mathrm{fin}}_{#1}}
\newcommand{\Ffrag}[2]{\mathcal{F}_{#1,#2}}
\newcommand{\D}{\mathcal{D}}
\newcommand{\out}{\mathrm{out}}
\newcommand{\tuple}[1]{\vec{#1}}
\newcommand{\blank}{\square}
\newcommand{\eps}{\varepsilon}
\newcommand{\pausesym}{\bot}
\newcommand{\fan}{\operatorname{fan}}
\newcommand{\YSub}[1]{\mathsf{YSub}(#1)}
\newcommand{\YSL}[2]{\mathsf{YSL}(#1,#2)}
\newcommand{\evalmap}{\mathsf{ev}}

\newcommand{\cutpos}[1]{%
  \raisebox{-0.7ex}{%
    \tikz[baseline=(char.base)]{%
      \node[draw,circle,inner sep=0.4pt,font=\scriptsize]
        (char) {\ensuremath{#1}};%
    }%
  }%
}

\title{Positive-Data Learning of Multiple Context-Free Languages under Fixed Finite-Monoid Typing}
\author{Takayuki Kuriyama\\
Independent Researcher, Tokyo, Japan\\
\texttt{growup.kuriyama@gmail.com}}
\date{}

\begin{document}
\maketitle

\begin{abstract}
Fix a fan-out bound $f$ and a homomorphism $h:\Sigma^*\to M$ into a finite
monoid.  We study positive-data learning of multiple context-free languages in
which tuples are compared only when corresponding components have equal
$h$-values, a finite-typed relaxation of Yoshinaka's multidimensional
substitutability.  We give a set-driven grammar reconstruction that becomes
exact after finitely many target-dependent positive observations and a
conservative learner that identifies every target in the limit; under a fixed
rule-rank bound, each update is polynomial in the accumulated sample.
For every fixed finite typing, rule-rank-one targets admit characteristic data
of total length polynomial in grammar size.  At arbitrary fixed rank, a
polynomial bound in grammar size and rule thickness is obtained whenever the
typing factors through a fixed prefix--suffix summary.  At every fan-out level
$m\ge2$, the $(1,1)$ boundary typing yields a rank-one target of minimum
fan-out $m$ outside Yoshinaka's class and outside the same-typing regular-filter
fragment.  At fan-out two, a separate finitely typed target lies outside
Yoshinaka's class even after arbitrary target-dependent regular filtering.
The thickness-sensitive bounds are stated for good presentations and transfer
polynomially from reduced $\lambda$-free presentations, even when deleting
rules are present; empty tuple components can cause exponential blow-up along
the standard normalization route.  For general higher-rank typings, an
original-thickness characteristic-data bound remains open.
\end{abstract}

\medskip
\noindent\textbf{Keywords:} grammatical inference; multiple context-free grammars; positive data; finite-monoid typing; multidimensional substitutability; Gold identification; rule rank.
\medskip

\section{Introduction}
\label{sec:introduction}

Positive-data grammatical inference asks when a language can be recovered from
an enumeration of its positive examples alone \cite{Gold1967,Angluin1980}.
Distributional approaches make this possible for restricted grammar classes by
turning observed interchangeability of substrings or tuples into grammar
nonterminals.  For context-free grammars (CFGs) this includes substitutable languages
\cite{ClarkEyraud2007} and Yoshinaka's $(k,\ell)$ hierarchy
\cite{Yoshinaka2008}.  Yoshinaka extended the same idea to multiple
context-free grammars (MCFGs), where nonterminals generate tuples and
substitution is tested by multi-hole sentence contexts \cite{Yoshinaka2009,Yoshinaka2011}.
MCFGs form a standard mildly context-sensitive formalism and can directly
represent discontinuous and cross-serial dependencies
\cite{SekiEtAl1991,Kallmeyer2010}.

This paper fixes, as part of the learnable class, a homomorphism
$h:\Sigma^*\to M$ into a finite monoid.  Two observed tuples are compared only
when their corresponding components have the same $h$-values.  The typing is
compositional because $h(uv)=h(u)h(v)$, and it is fixed before the positive
text is seen.  This separates the problem of choosing a useful finite bias from
the problem of learning under that bias, in the same spirit as typing or domain
bias in automata inference \cite{CosteTyping2004}.  Prefix--suffix morphisms
provide one systematic family of such biases, while transition morphisms of
finite automata provide another.  The fan-out bound is part of the same bias:
it controls not only representational power but also which tuple arities are
tested for substitution, so choosing a larger $f$ is not automatically safer.
If the typing itself is unknown but a bound on the size of its finite codomain
is fixed, taking the product of all homomorphisms into monoids up to that size
reduces the bounded union to one fixed typing.  Without a
codomain-size bound, the union contains all regular languages and is not
TxtEx-identifiable from positive data.  A useful way to read the fixed typing
is as bounded finite-state information that selects among finitely many
structural regimes while the MCFG carries the unbounded dependency itself.
The framework does not attempt to encode an unbounded copied feature sequence;
the separation examples later make both the usefulness and this limitation
explicit.

Other learning biases organize the inference problem differently.  Takada's
control-set approach represents even linear context-free languages by a fixed grammar whose
derivations are selected by a regular control language
\cite{Takada1988,Takada1995}; here the fixed object instead acts on terminal
yields and gates distributional substitutions.  Coste, Garet, and Nicolas
localize the evidence used to infer substitutability: bounded immediate
neighborhoods may license a comparison even when the outer sentence contexts
differ \cite{CosteGaretNicolas2012}; see also the polynomial reduction
approach of Coste and Nicolas~\cite{CosteNicolas2019}.  Here the evidence remains a common full
sentence context, while the finite typing records a compositional summary of
the compared factor or tuple component itself.  For MCFGs and related
formalisms, richer-information settings include learning from positive data
together with membership queries \cite{Yoshinaka2010} and distributional
learning of parallel MCFGs \cite{ClarkYoshinaka2014}.  These comparisons
motivate the fixed-typing bias but are not used in the proofs.

\paragraph{Contributions.}
The paper has three main contributions.
\begin{enumerate}[label=\textup{(\arabic*)},leftmargin=2.4em]
\item \emph{Exact reconstruction and learning under a fixed finite typing.}
For every fixed fan-out bound and finite typing, observed tuple values,
sample-local MCFG compositions, and equal-type links yield a set-driven
reconstruction that is exact after a finite target-dependent positive sample.
No rule-rank cap is needed for this qualitative theorem.  A conservative
sequential wrapper gives TxtEx identification; when rule rank is fixed, every
reconstruction step is polynomial in the accumulated positive sample.
\item \emph{Characteristic-data bounds.}
For arbitrary fixed $h$, rule-rank-one targets have characteristic data of
total length polynomial in grammar size.  At arbitrary fixed rank, typings
that factor through a fixed prefix--suffix summary admit a polynomial bound in
grammar size and rule thickness.  The proof simultaneously preserves bounded
boundary information in several tuple components.  Two finite refinements can
obstruct a direct transfer from an arbitrary presentation: full yield-type
refinement can increase rule thickness exponentially, and Appendix~\ref{app:normalization-details}
shows the same phenomenon for empty-component elimination.
\item \emph{Multiple-context expressiveness and fan-out sensitivity.}
For every $m\ge2$, a rank-one language of minimum fan-out $m$ belongs to the
$(1,1)$-typed class while lying outside Yoshinaka's class and the regular-filter
fragment recognized by that same typing.  A separate fan-out-two example lies
outside Yoshinaka's class followed by \emph{any} target-dependent regular
filter.  Increasing the tested tuple arity can already destroy substitutability
in Yoshinaka's untyped baseline; our stronger fan-out-two witness shows that
this loss need not be repairable by \emph{any} finite typing.
\end{enumerate}

\paragraph{Relation to prior work and scope.}
With trivial typing, the semantic baseline is Yoshinaka's multidimensional
substitutability~\cite{Yoshinaka2011}.  Yoshinaka also removes the rule-rank
bound for qualitative identification~\cite[Corollary~2]{Yoshinaka2011}; the
qualitative contribution here extends that rank-unbounded setting from trivial
typing to a fixed finite typing and uses a sample-local tuple constructor whose
relation to Yoshinaka's finite-sample grammar is made explicit in
Section~\ref{sec:learner}.  At fan-out one, the semantic target class is exactly
the fixed-$h$ CFG class of the companion paper~\cite{kuriyamaCFG}.  The
qualitative theorem, the rank-one/linear-CFG data bound, and the fixed-window
compatibility are the specific fan-out-one specializations already established
there; Proposition~\ref{prop:fanout-one-operator} below relates the two
finite-sample reconstruction operators themselves.  Polynomial-time statements
below treat the displayed fan-out, rank cap, and finite typing as fixed.
Characteristic-data bounds concern the set-driven reconstruction; the
conservative learner is used separately for syntactic stabilization.  The
higher-rank thickness obstructions concern the displayed normalization and
refinement routes, not lower bounds on all possible learners or characteristic
samples.

Section~\ref{sec:prelim} fixes notation and the learning model.
Section~\ref{sec:learner} gives the reconstruction and convergence proofs.
Section~\ref{sec:char-data} develops the quantitative bounds and their
limitations.  Section~\ref{sec:comparison} compares the framework with
Yoshinaka's classes and gives the separation examples.

\section{Preliminaries}
\label{sec:prelim}

Throughout, $\Sigma$ denotes a finite terminal alphabet, $\eps$ denotes the
empty word, $\Sigma^*$ is the set of all finite words over $\Sigma$, and
$\Sigma^+:=\Sigma^*\setminus\{\eps\}$.

\subsection{Positive-Data Learning}

We use Gold's explanatory learning model and restate the needed conventions
so that the MCFG development is self-contained.

\begin{definition}[Text and TxtEx identification]
We use Gold's explanatory identification criterion~\cite{Gold1967}.  To handle the empty language, fix a pause symbol
\(\pausesym\notin\Sigma\); thus
\(\pausesym,\pausesym,\ldots\) is a text for the empty language.
A \emph{text} for \(L\subseteq\Sigma^*\) is an infinite sequence over
\(L\cup\{\pausesym\}\) in which every element of \(L\) appears at least once.
Pause symbols may occur arbitrarily, and \(\pausesym\) is not an end-of-text marker.

The \emph{content} of a finite prefix is the finite set
\(K\subseteq L\) of positive examples from \(\Sigma^*\) that occur in that
prefix.  Thus neither the positions nor the number of pause symbols, nor repeated
occurrences of the same positive example, affect the content.
A learner is a computable function mapping each finite prefix to a hypothesis
grammar.

A learner \emph{TxtEx-identifies} \(L\) from positive data if, for every text
for \(L\), there are \(n_0\) and a single hypothesis grammar \(H\) such that
all stages from \(n_0\) onward output exactly \(H\), with \(L(H)=L\).  It
TxtEx-identifies a class when it does so for every language in the class.  A
sequential learner is \emph{conservative} if it keeps its current hypothesis
whenever the newly observed positive datum is already in that hypothesis
language, changing it only when the datum is not generated.
\end{definition}

For later use, fix a distinguished hypothesis code $H_0$ with
$L(H_0)=\emptyset$.

\begin{definition}[Set-driven reconstruction operator and characteristic sample]
\label{def:characteristic-sample}
A \emph{set-driven reconstruction operator} \(\mathcal B\) maps each finite
positive sample \(K\) to a grammar \(\mathcal B(K)\), depending only on the
set \(K\), not on presentation order or multiplicity.  It is
\emph{sample-consistent} if \(K\subseteq L(\mathcal B(K))\) for every
finite \(K\).

Following the characteristic-set viewpoint of de~la~Higuera
\cite{deLaHiguera1997}, for a target language \(L\), a finite set
\(K^\star\subseteq L\) is a \emph{characteristic sample for \(\mathcal B\)}
if, for every finite set
\(K\), \(K^\star\subseteq K\subseteq L\) implies
\(L(\mathcal B(K))=L\).  The case \(K^\star=\emptyset\) is allowed.  We
use \emph{characteristic data} as a neutral shorthand for the positive content
of such a sample; whenever a polynomial bound is claimed, we state explicitly
whether it concerns cardinality or total positive length.

This finite-sample notion is distinct from TxtEx learning: \(\mathcal B(K)\)
need not stabilize syntactically as the accumulated sample grows.  The
sequential learner will be introduced separately after exact reconstruction is
proved.
\end{definition}

\begin{definition}[Positive sample size]
For finite \(K\subseteq\Sigma^*\), define
\(\|K\|_+:=\sum_{w\in K}(|w|+1)\).
For a tuple \(\tuple u=(u_1,\ldots,u_d)\), write
\(\|\tuple u\|_1:=\sum_{i=1}^d |u_i|\)
for its total component length.
\end{definition}


\subsection{Finite-Monoid Typing}

\begin{definition}[Finite typing representation]
\label{def-monoid-morphism}
A finite-monoid homomorphism \(h\colon\Sigma^*\to M\) is represented by the
multiplication table of \(M\) together with the letter images \(h(a)\),
\(a\in\Sigma\).  Hence \(h(w)\) is computable by a left-to-right product in
linear time in \(|w|\).  For a tuple \(\tuple w=(w_1,\ldots,w_d)\), write
\(h^{(d)}(\tuple w):=(h(w_1),\ldots,h(w_d))\in M^d\).
\end{definition}

\begin{definition}[Refinement of finite typings]
\label{def:h-refinement}
Let \(h\colon\Sigma^*\to M\) and \(h'\colon\Sigma^*\to M'\) be finite
homomorphisms.  We say that \(h'\) \emph{refines} \(h\), and write
\(h\preceq h'\), when
\[
  \ker h'\subseteq\ker h.
\]
Equivalently, the value of \(h\) factors through the image of \(h'\): there is
a unique monoid homomorphism
\(\upsilon\colon\operatorname{im}h'\to\operatorname{im}h\) satisfying
\(h=\upsilon\circ h'\), after restricting the codomains to their images.
\end{definition}

For two typings $h$ and $g$, their product typing $(h,g)$ refines both and
satisfies $\ker(h,g)=\ker h\cap\ker g$.

Finite-monoid homomorphisms serve here as finite-state typings.  Let
$h_{\mathrm{triv}}:\Sigma^*\to\{1\}$ denote the unique homomorphism into the
one-element monoid.
For example, the transition morphism of a deterministic finite automaton (DFA) for a regular language
$L_{\mathrm{reg}}\supseteq L$ supplies such an $h:\Sigma^*\to M$, where $h(w)$ records the
state transformation induced by $w$ \cite{Pin1986}.  The reconstruction uses
only this finite summary, not $L_{\mathrm{reg}}$ as a membership oracle for $L$.


We call $h$ a \emph{finite-monoid typing} and $h^{(d)}(\tuple w)$ the
\emph{yield type} of the tuple.  At fan-out one the semantic target class is
the fixed-$h$ CFG class of~\cite{kuriyamaCFG}; the relation between the two
reconstruction operators is stated in Proposition~\ref{prop:fanout-one-operator}.
The prefix--suffix morphism $h_\Sigma^{k,\ell}$ below is the same boundary
typing denoted $h_{k,\ell}$ there.

\begin{definition}[Prefix--suffix typing]
\label{def:prefix-suffix-typing}
For $k,\ell\ge0$, put $m_0:=\max\{1,k+\ell\}$.  For a word $w$, let
$\operatorname{pref}_k(w)$ and $\operatorname{suf}_\ell(w)$ denote its prefix
and suffix of the indicated lengths whenever $|w|\ge m_0$, with the length-zero
windows equal to $\eps$.  Define
\[
 h_\Sigma^{k,\ell}(w)=
 \begin{cases}
   w, & |w|<m_0,\\
   \langle\operatorname{pref}_k(w),\operatorname{suf}_\ell(w)\rangle,
      & |w|\ge m_0,
 \end{cases}
\]
using disjoint tags for the short-word and prefix--suffix elements.  Put
$M_{k,\ell}:=h_\Sigma^{k,\ell}(\Sigma^*)$.  For short words $x,y$ and long
elements $\langle u,v\rangle$, define
\[
\begin{aligned}
 x\cdot y&:=h_\Sigma^{k,\ell}(xy), &
 x\cdot\langle u,v\rangle&:=\langle\operatorname{pref}_k(xu),v\rangle,\\
 \langle u,v\rangle\cdot y&:=\langle u,\operatorname{suf}_\ell(vy)\rangle, &
 \langle u,v\rangle\cdot\langle u',v'\rangle&:=\langle u,v'\rangle.
\end{aligned}
\]
\end{definition}

With this multiplication,
$h_\Sigma^{k,\ell}(xy)=h_\Sigma^{k,\ell}(x)\cdot h_\Sigma^{k,\ell}(y)$.
Since $h_\Sigma^{k,\ell}$ is surjective, associativity is inherited from
concatenation and $h_\Sigma^{k,\ell}(\eps)$ is the identity.  Thus
$h_\Sigma^{k,\ell}:\Sigma^*\to M_{k,\ell}$ is a finite-monoid homomorphism.

A finite typing $h$ is called \emph{$(k,\ell)$-boundary-determined} when
$h\preceq h_\Sigma^{k,\ell}$.  Equivalently, equality of prefix--suffix types implies equality of $h$-types;
a factor map exists on the image of the prefix--suffix typing.


\subsection{Multiple Context-Free Grammars}

We use the standard tuple-generating definition of multiple context-free
grammars \cite{SekiEtAl1991}; see also \cite{Kallmeyer2010} for a general
overview of MCFGs and related mildly context-sensitive formalisms.
The general-rank notation below is used throughout the reconstruction and proofs.

\begin{definition}[Standard MCFG presentation]
\label{def:standard-mcfg}
A \emph{multiple context-free grammar} is a tuple
\(G=(V,\Sigma,P,S,\fan)\), where \(V\) is a finite set of nonterminals,
\(\Sigma\) is a finite terminal alphabet, \(S\in V\) has \(\fan(S)=1\), and
\(\fan(A)\ge1\) is the fan-out of each \(A\in V\).  A rule \(R\) of rank \(r\ge0\) has the form
\[
  R:\quad A\to\rho(B_1,\ldots,B_r),
  \qquad
  \rho=(\alpha_1,\ldots,\alpha_{d_0}),
\]
where \(\rho\) is the output template, \(d_0=\fan(A)\), \(d_j=\fan(B_j)\), and, with
\[
  \Xi_j:=\{\xi_j^1,\ldots,\xi_j^{d_j}\},
  \qquad
  \Gamma_\rho:=\Sigma\uplus\Xi_1\uplus\cdots\uplus\Xi_r,
\]
each \(\alpha_i\in\Gamma_\rho^*\) and every variable in
\(\Xi_1\uplus\cdots\uplus\Xi_r\) occurs at most once in the complete tuple
\((\alpha_1,\ldots,\alpha_{d_0})\).  We call this at-most-once condition
\emph{linearity}.  The rule is \emph{nondeleting} if every such variable occurs
exactly once.  Thus nondeleting does not mean that derived tuple components are
nonempty; empty strings remain allowed.

For child tuples
\(\tuple u_j=(u_{j,1},\ldots,u_{j,d_j})\in(\Sigma^*)^{d_j}\), simultaneous substitution of
\(u_{j,\kappa}\) for \(\xi_j^{\kappa}\) defines
\(\rho(\tuple u_1,\ldots,\tuple u_r)\in(\Sigma^*)^{d_0}.\)
When \(r=0\), the rule is a constant rule and derives the tuple
\((\alpha_1,\ldots,\alpha_{d_0})\).  The tuple languages \(L_A(G)\) form the least family, ordered componentwise by
inclusion, that is closed under all rule applications, and
\(L(G):=\{w\in\Sigma^*:(w)\in L_S(G)\}.\)
Equivalently, a finite derivation tree rooted at $A$ is evaluated bottom-up by
the rule templates.  We write $A\Rightarrow_G^*\tuple u$ when such a tree
yields $\tuple u\in(\Sigma^*)^{\fan(A)}$; this is a \emph{terminal
derivation}, and $\tuple u$ is its \emph{terminal tuple}.  A terminal
derivation $S\Rightarrow_G^*(w)$ from the start symbol is called
\emph{successful}.
An \(f\)-MCFG has \(\fan(A)\le f\) for every nonterminal, and an \(f\)-MCFL is
a language generated by an \(f\)-MCFG.  A grammar has \emph{rule-rank cap} \(r_{\max}\) if every rule has rank at most \(r_{\max}\).
\end{definition}

At fan-out one, the formalism includes ordinary context-free grammars.  A
\emph{linear context-free grammar} is a CFG in which every production has at
most one nonterminal on its right-hand side; a \emph{linear context-free
language} is generated by such a grammar.  This is distinct from the
``linearity'' condition on MCFG rule templates above, which means that each
rule variable occurs at most once.

Fix a standard explicit symbol-level encoding of finite MCFG presentations and
write $\|G\|$ for its length.  The encoding counts the nonterminal and rule
records and every terminal or variable occurrence in each rule template.  All
polynomial grammar-size bounds below refer to this encoding.

For integers \(f\ge1\) and \(r\ge0\), write
\[
\begin{aligned}
  \mathsf{MCFL}(f,r):=\{L(G):{}&G\text{ is an MCFG of fan-out at most }f,\\
                   &\text{with rule rank at most }r\}.
\end{aligned}
\]
Thus \(f\text{-}\mathrm{MCFL}=\bigcup_{r<\infty}\mathsf{MCFL}(f,r)\), because every
finite MCFG presentation has finite rule rank.
When minimum rule rank at fan-out \(f\) is mentioned below, it means the
language-level minimum over all fan-out-\(f\) presentations, not the rank of a
particular supplied grammar.


We use $f$ throughout for fan-out/dimension and $r$ for a generic rule rank;
$r_{\max}$ denotes a prescribed rank cap.  Rule-template variables are always
written $\xi_j^\kappa$, so that $x,y,z$ (with or without arrows) remain free for
words and observed tuples.  Terminals in examples, including digits and
$\mathtt\#$, are typeset in \texttt{typewriter} font.


When two nonterminals have the same fan-out \(d\), we use
\(A\to B\)
as shorthand for the unary identity-template rule
\(A\to(\xi_1^1,\ldots,\xi_1^d)(B)\).  This denotes only the identity
correspondence between components: at fan-out two, for example, it means
\(A\to(\xi_1^1,\xi_1^2)(B)\), not a swapping rule such as
\(A\to(\xi_1^2,\xi_1^1)(B)\).
Likewise, \(A\to\tuple c\) denotes a rank-zero
constant rule.  These abbreviations are used for typed start links and for Rule~(Link) of the
reconstruction below.

The standard formalism allows rank zero, nonidentity unary rules, arbitrary
finite rule rank, and deleting rules.  For a nondeleting rule
\(R:A\to\rho(B_1,\ldots,B_r)\), with
\(\rho=(\alpha_1,\ldots,\alpha_{d_0})\), is
\emph{nonpermuting} if, for each child \(B_j\), the variables
\(\xi_j^1,\ldots,\xi_j^{\fan(B_j)}\) occur in this order when the parent components
are scanned from left to right.  Following Yoshinaka
\cite[Section~2.3]{Yoshinaka2011}, a rule is \emph{$\lambda$-free} if no output
template $\alpha_i$ is empty, and a nonpermuting rule is \emph{nonmerging} if
no $\alpha_i$ contains adjacent variables $\xi_j^{\kappa} \xi_j^{\kappa+1}$ of the same child.
A presentation is \emph{good} if all its rules are $\lambda$-free,
nondeleting, nonpermuting, and nonmerging.  We write $\eps$ for the empty word
throughout; ``$\lambda$-free'' is retained only as Yoshinaka's standard name
for this rule condition.  We use good presentations in the completeness proof.  The nonmerging condition ensures that, after the siblings
of a child are fixed, consecutive components of that child remain separated
by nonempty terminal material or by a nonempty sibling yield.  This matches
Yoshinaka's multidimensional-context convention.  In particular,
$\lambda$-freeness implies by induction that every component of every terminal tuple derived from a non-start
nonterminal is nonempty.

The normalization used below combines the standard rank-preserving nondeleting
conversion of Seki et al. and Kanazawa with component permutation, explicit
empty-component removal, and Yoshinaka's nonmerging conversion
\cite[proof of Lemma~2.2]{SekiEtAl1991}\cite[Section~2.1]{Kanazawa2019Ogden}\cite[Section~2.3, Lemma~1]{Yoshinaka2011}.  Yoshinaka also notes immediately after that
lemma that, for fixed dimension and rule rank, these standard conversions yield
an equivalent good grammar of size $O(\|G\|)$~\cite[Section~2.3]{Yoshinaka2011}.

Fan-out and rule rank are distinct grammar resources.  In particular, there is
no blanket binarization theorem that preserves an arbitrary fixed fan-out for
general MCFG/LCFRS presentations: reducing rule rank may require increased
fan-out \cite{RambowSatta1999,GomezRodriguezEtAl2009}.
Accordingly the qualitative reconstruction below is defined without imposing a
binary normal form or a known rank cap.


\paragraph{Reduced grammars.}
The nonterminal dependency graph of an MCFG has an edge $A\to B$ whenever
$B$ occurs as a child in a rule with left-hand side $A$.
A nonterminal is \emph{reachable} if it is reachable from the start symbol in
this graph, and \emph{productive} if its tuple language is nonempty.
The grammar is \emph{reduced} if every nonterminal is both reachable and
productive.


\subsection{Normalization}

\begin{proposition}[Good normalization with rank preservation]
\label{prop:good-normalization}
Let $G_{\mathrm{in}}$ be a finite standard $f$-MCFG of maximum rule rank at
most $r_{\max}$.  There is a finite good $f$-MCFG $G_{\mathrm{good}}$ of
maximum rule rank at most $r_{\max}$ such that
\[
  L(G_{\mathrm{good}})=L(G_{\mathrm{in}})\setminus\{\eps\}.
\]
For fixed $f$ and $r_{\max}$,
$\|G_{\mathrm{good}}\|=O_{f,r_{\max}}(\|G_{\mathrm{in}}\|)$.
If $\eps\notin L(G_{\mathrm{in}})$, the two grammars are equivalent.
\end{proposition}

\begin{proof}
First delete unreachable or unproductive nonterminals, then apply the
rank-preserving nondeleting normalization of Seki et al., in the form recorded by
Kanazawa~\cite[proof of Lemma~2.2]{SekiEtAl1991}\cite[Section~2.1]{Kanazawa2019Ogden}.
Decorating nonterminals by component permutations makes the grammar
nonpermuting.  Next decorate by nonempty-component patterns realizable by terminal
derivations, as made explicit in Appendix~\ref{app:normalization-details}, and
project away empty components; this removes $\eps$ at the start and makes every
retained output template $\lambda$-free.  Finally apply Yoshinaka's nonmerging
conversion~\cite[Section~2.3, Lemma~1]{Yoshinaka2011}.
None of these steps increases fan-out or rule rank.  With $f,r_{\max}$ fixed,
component-pattern choices contribute at most $2^f$ nonterminal decorations and
$2^{f r_{\max}}$ child-pattern choices per rule, while permutation decoration
contributes at most $f!$ copies.  Yoshinaka's conversion has the same
fixed-parameter linear-size behavior.  Hence the total encoding size remains
$O_{f,r_{\max}}(\|G_{\mathrm{in}}\|)$.  The derivation correspondence and the
empty-component construction are recorded in
Appendix~\ref{app:normalization-details}.
\end{proof}

This proposition controls size, fan-out, and rule rank, but it does not identify
the rule thickness of an arbitrary deleting presentation with the thickness
after that nondeleting normalization.  Thickness-sensitive statements are
therefore first proved for good presentations.  For already $\lambda$-free,
nondeleting presentations, the structural permutation and nonmerging steps do
admit a polynomial thickness transfer; Proposition~\ref{prop:structural-good-thickness}
records it explicitly.

\subsection{Tuple Contexts and Substitutability}

\paragraph{Sentence contexts.}
A \emph{sentence context of arity $d$} is a word
$E=v_0\blank_1v_1\cdots v_{d-1}\blank_dv_d$ in which each named hole occurs
exactly once, $v_0,v_d\in\Sigma^*$, and every internal separator
$v_1,\ldots,v_{d-1}$ belongs to $\Sigma^+$.  Let $\mathcal E_d$ be the set
of such contexts.  This is Yoshinaka's multidimensional-context convention
\cite[Section~3.1]{Yoshinaka2011}.
For $\tuple w=(w_1,\ldots,w_d)$, write $E[\tuple w]$ for the result of filling
$\blank_i$ by $w_i$, and put
$|E|_{\mathrm{fix}}:=\sum_{i=0}^d|v_i|$.  For brevity, we say that a
sentence context $E$ \emph{accepts} $\tuple w$ in a language $L$ when
$E[\tuple w]\in L$.

\begin{definition}[Tuple occurrence in a sample word]
\label{def:tuple-occurrence}
Let $K\subseteq\Sigma^*$ and $w=a_1\cdots a_n\in K$.  Number the cut
positions of $w$ as
$\cutpos{0}a_1\cutpos{1}\cdots a_n\cutpos{n}$.
For $0\le i<i'\le n$, write
$w[i:i']:=a_{i+1}\cdots a_{i'}$.

An \emph{arity-$d$ tuple occurrence} in $w$ is a sequence of nonempty,
left-to-right intervals
\[
  0\le i_1<i_1'<i_2<i_2'<\cdots<i_d<i_d'\le n.
\]
It determines the tuple
$\tuple x=(w[i_1:i_1'],\ldots,w[i_d:i_d'])\in(\Sigma^+)^d$ and the
context $E\in\mathcal E_d$ obtained by replacing those intervals by
$\blank_1,\ldots,\blank_d$.  We identify the occurrence with $(E,\tuple x)$.
For example, the two intervals spelling $\mathtt a$ and $\mathtt c$ in $\mathtt{abc}$ give the tuple
$(\mathtt a,\mathtt c)$ and the context $\blank_1\mathtt b\blank_2$.
\end{definition}

For $L\subseteq\Sigma^*$ and $\tuple x\in(\Sigma^*)^d$, define its
\emph{$d$-tuple distribution}
\[
  \D_L^{(d)}(\tuple x):=\{E\in\mathcal E_d:E[\tuple x]\in L\}.
\]
When $d$ is clear, the superscript is omitted.

\begin{definition}[$(f,h)$-tuple substitutability]
\label{def:tuple-substitutable}
Let $f\ge1$ and let $h:\Sigma^*\to M$ be a finite monoid
homomorphism.  A language $L\subseteq\Sigma^*$ is
\emph{$(f,h)$-tuple-substitutable} if, for every $1\le d\le f$ and all
$\tuple x,\tuple y\in(\Sigma^+)^d$,
\[
 h^{(d)}(\tuple x)=h^{(d)}(\tuple y),\qquad
 \D_L^{(d)}(\tuple x)\cap\D_L^{(d)}(\tuple y)\ne\emptyset
\]
imply $\D_L^{(d)}(\tuple x)=\D_L^{(d)}(\tuple y)$.
Thus empty strings are not compared as tuple components; when $\eps$ belongs
to the target, it is handled separately at the start symbol.  This is Yoshinaka's nonempty-component convention for multidimensional
substitutability~\cite{Yoshinaka2011}; see also Clark~\cite{Clark2013}.  The restriction is
substantive already at fan-out one: if $\eps$ itself were allowed among the
compared fragments, the ordinarily substitutable language $\mathtt a^+$ would fail
the condition because $\mathtt a$ and $\eps$ share a context but do not have the same
distribution.
\end{definition}

At fan-out one, every sentence context has the form
$u\blank_1v$, so $\D_L^{(1)}(x)$ is canonically the ordinary two-sided
distribution $\D_L(x)=\{(u,v):uxv\in L\}$.  Hence $(1,h)$-tuple
substitutability is precisely the fixed-$h$ CFG condition of
\cite{kuriyamaCFG}.

\paragraph{Typing--distribution equivalence.}
For $\tuple x,\tuple y\in(\Sigma^+)^d$, write
$\tuple x\equiv_{L,h}^d\tuple y$ if and only if
\[
 h^{(d)}(\tuple x)=h^{(d)}(\tuple y)
 \qquad\text{and}\qquad
 \D_L^{(d)}(\tuple x)=\D_L^{(d)}(\tuple y).
\]
For fixed $L,h,d$ this is an equivalence relation; the superscript is omitted
when the arity is clear.

\begin{definition}[The target class]
\label{def:class-cmcf}
Fix $f\ge1$ and a finite monoid homomorphism $h$.  The class
$\Cmcf{f}{h}$ consists of all $f$-MCFLs that are
$(f,h)$-tuple-substitutable.
\end{definition}

If $h\preceq h'$, then
$\Cmcf{f}{h}\subseteq\Cmcf{f}{h'}$: equality of the finer $h'$-types implies
equality of the coarser $h$-types, so the substitution conclusion available
under $h$ applies.  At fan-out one this is exactly the fixed-$h$ CFG class
$\mathcal C_h^{\mathrm{cf}}$ of~\cite{kuriyamaCFG}.

For every $f\ge1$,
\[
  \Cmcf{f+1}{h}\cap f\text{-}\mathrm{MCFL}
  \subseteq \Cmcf{f}{h},
\]
because the fan-out-$f$ condition tests only the arities
$1,\ldots,f$.  The possible loss of substitutability when the tested arity is
increased is already present for trivial typing: Yoshinaka's
$L_{\mathrm{reverse}}=\{w\#w^R\}$ is one-dimensional substitutable but not
two-dimensional substitutable~\cite[Example~5]{Yoshinaka2011}.  What finite
typing adds here is the stronger phenomenon of
Proposition~\ref{prop:strong-fanout-separation}: a target can be learnable at
fan-out one under a finite typing yet fail the fan-out-two condition for
\emph{every} finite typing.  Thus $f$ is part of the learning bias, not merely
an upper bound on representational power.  No monotonicity of the learned
classes in $f$ is claimed.

For every $\mathsf{Acc}\subseteq M$, the language
$h^{-1}(\mathsf{Acc})$ belongs to $\Cmcf{f}{h}$.  Indeed, it is regular, and equal componentwise $h$-types give the
same total $h$-value in every sentence context, hence equal tuple
distributions.

\begin{proposition}[Intersection and abelian filtering]
\label{prop:abelian-coset-filter}
Fix $f$ and $h$.
\begin{enumerate}[label=(\roman*),leftmargin=*]
\item If $L_1$ and $L_2$ are $(f,h)$-tuple-substitutable, then
$L_1\cap L_2$ is $(f,h)$-tuple-substitutable.
\item Let $L$ be $(f,h)$-tuple-substitutable, let $\mathbb A$ be an abelian
group, let $\varphi:\Sigma^*\to\mathbb A$ be a monoid homomorphism, let
$\mathcal H\le\mathbb A$, and let $\gamma\in\mathbb A$.  Then
$L\cap\varphi^{-1}(\gamma+\mathcal H)$ is $(f,h)$-tuple-substitutable.
\end{enumerate}
Whenever the resulting intersection is an $f$-MCFL, it therefore belongs to
$\Cmcf{f}{h}$.
\end{proposition}

\begin{proof}
For (i), equal-type tuples $\tuple x,\tuple y$ with some context $E$ satisfying
$E[\tuple x],E[\tuple y]\in L_1\cap L_2$ have equal tuple distributions in each $L_i$, hence in their intersection.
For (ii), first observe that $\varphi^{-1}(\gamma+\mathcal H)$ is
$(f,h_{\mathrm{triv}})$-tuple-substitutable.  Indeed, if a common context
accepts tuples $\tuple x,\tuple y$, then
$\sum_i(\varphi(y_i)-\varphi(x_i))\in\mathcal H$.  By commutativity, for every other sentence
context $E'$ the difference
$\varphi(E'[\tuple y])-\varphi(E'[\tuple x])$ is the same element of
$\mathcal H$, so coset membership is preserved in both directions.  Trivial
typing is coarser than $h$, hence the same coset language is
$(f,h)$-tuple-substitutable.  Apply (i).
\end{proof}

As an immediate consequence, for any $\mathsf{Acc}\subseteq M$, abelian group
$\mathbb A$, homomorphism $\varphi:\Sigma^*\to\mathbb A$, subgroup
$\mathcal H\le\mathbb A$, and $\gamma\in\mathbb A$, every $f$-MCFL of the
form
\[
 h^{-1}(\mathsf{Acc})\cap\varphi^{-1}(\gamma+\mathcal H)
\]
belongs to $\Cmcf{f}{h}$.  Thus one typing can support a family of targets
while the abelian constraint varies; these factorizations are not supplied to
the learner.

\section{Exact Reconstruction from Positive Data}
\label{sec:learner}

We define the finite-sample reconstruction operator, prove soundness, and then
prove completeness using a normalized target grammar refined by yield type.
The target grammar is used only to select a finite characteristic sample.
Since every observed tuple component is nonempty, the rank of an enumerated
composition witness is bounded by the length of the sample word containing it.

\subsection{Set-Driven Reconstruction and Sample-Bounded Rank}


Following Definition~\ref{def:tuple-occurrence}, enumerate all tuple occurrences
of arity at most $f$ in the sample $K$.  The reconstruction uses one nonterminal
$[\tuple x]$ for every observed tuple value $\tuple x\in(\Sigma^+)^d$, together
with a fresh start symbol $\widehat S$.  It is not told which occurrences
correspond to target nonterminals or rules; all admissible segmentations are
enumerated.

\begin{definition}[Composition witness]
\label{def:composition-witness}
Let $r\ge1$.  A \emph{composition witness of rank $r$} in $K$ consists of an
observed parent occurrence $(E_{\mathrm{par}},\tuple z)$ of arity $e\le f$ together with a
segmentation of the parent components into terminal gaps and nonempty labelled
intervals.  The labels are variables $\xi_j^{\kappa}$ with
$1\le j\le r$ and $1\le\kappa\le d_j\le f$; every variable occurs exactly once,
and for each fixed $j$ the variables
$\xi_j^1,\ldots,\xi_j^{d_j}$ occur in this order when the parent components are
scanned from left to right.  No output component may contain adjacent variables
$\xi_j^{\kappa}\xi_j^{\kappa+1}$ of the same child.  The labelled intervals
induce child tuples $\tuple u_j=(u_{j,1},\ldots,u_{j,d_j})$ and determine a
unique linear template $\rho$ satisfying
\[
  \tuple z=\rho(\tuple u_1,\ldots,\tuple u_r).
\]
\end{definition}

\begin{lemma}[Materialization of child tuples]
\label{lem:composition-witness-materialization}
Every child tuple induced by a composition witness is itself an observed tuple
occurrence, and the induced template $\rho$ is good.
\end{lemma}

\begin{proof}
Because the parent is an observed tuple occurrence, the components of each
$\tuple u_j$ remain separated in the same sample word.  Components lying in one
parent component are separated by the terminal material forced by the
nonmerging condition, while components lying in different parent components
are separated by the nonempty internal material of the parent sentence
context.  Thus each $\tuple u_j$ is an observed tuple occurrence.  Every
variable occurs exactly once, the components of each child occur in order, and
the definition excludes adjacent consecutive components of one child; hence
$\rho$ is linear, nondeleting, nonpermuting, and nonmerging.  Finally every
parent component is nonempty, so every output component of $\rho$ is nonempty.
Thus $\rho$ is good.
\end{proof}

\begin{lemma}[Enumeration and sample-bounded rank]
\label{lem:general-witness-enumeration}
Fix $f$ and a finite sample $K$, and put $n_K:=\max\{1,\|K\|_+\}$.  If a
rank-$r$ composition witness has parent tuple $\tuple z$ occurring in a sample
word $w$, then
\[
  r\le\|\tuple z\|_1\le |w|\le\|K\|_+.
\]
All composition witnesses are effectively enumerable.  Under a rule-rank cap
$r_{\max}$, those of rank at most $r_{\max}$ are enumerable in
$n_K^{O(f r_{\max})}$ time and space.
\end{lemma}

\begin{proof}
Every child tuple has at least one nonempty component.  Linearity and
nondeletion make the child components contribute disjointly to the parent, so
$\|\tuple z\|_1\ge\sum_j\|\tuple u_j\|_1\ge r$; the remaining inequalities
follow because the parent components are disjoint intervals of $w$.

An arity-$d$ tuple occurrence is determined by $2d$ cuts.  For a fixed rank
$r$, a candidate segmentation has at most $fr$ labelled child intervals, whose endpoints are chosen among the $|w|+1\le n_K+1$ cut positions of
the parent sample word.  For $r\le r_{\max}$, assignments to parent components and
terminal gaps contribute only a factor depending on $f$ and $r_{\max}$, and
the nonpermuting and nonmerging conditions are directly checkable.  Summing
over $r\le r_{\max}$ gives the stated bound.  Without a cap, the displayed
rank inequality restricts the effective enumeration to
$1\le r\le\|K\|_+$.
\end{proof}

\begin{definition}[Capped and uncapped reconstruction operators]
\label{def:reconstruction}
Fix $f$ and $h$.  For $r_{\max}\ge1$, let
$\mathcal B_{f,h}^{\le r_{\max}}(K)$ be the grammar
$\widehat G_{\le r_{\max}}(K)$ with a fresh start symbol $\widehat S$, one nonterminal $[\tuple x]$ for each nonempty tuple value $\tuple x$ observed in $K$,
and the following rules:
\begin{enumerate}[leftmargin=*]
\item[\textnormal{(Start)}] for every $w\in K\cap\Sigma^+$,
      $\widehat S\to[(w)]$, and if $\eps\in K$, also
      $\widehat S\to\eps$;
\item[\textnormal{(Const)}] for every observed tuple, the constant rule
      $[\tuple u]\to\tuple u$;
\item[\textnormal{(Comp)}] for every composition witness of rank
      $1\le r\le r_{\max}$ with parent tuple $\tuple z$, template $\rho$,
      and induced child tuples $\tuple u_1,\ldots,\tuple u_r$, the rule
      \[
        [\tuple z]\to\rho([\tuple u_1],\ldots,[\tuple u_r]);
      \]
\item[\textnormal{(Link)}] for observed nonempty tuples
      $\tuple x,\tuple y$ of the same arity $d$, the identity-template rule
      $[\tuple x]\to[\tuple y]$ whenever
      $h^{(d)}(\tuple x)=h^{(d)}(\tuple y)$ and some $E\in\mathcal E_d$
      satisfies $E[\tuple x],E[\tuple y]\in K$.
\end{enumerate}
If $K\ne\emptyset$, define
$\mathcal B_{f,h}(K):=\widehat G(K):=
\widehat G_{\le\|K\|_+}(K)$; for $K=\emptyset$, let
$\mathcal B_{f,h}(\emptyset)$ be the grammar with no productive start rule.
Fixed encoding orders make both reconstruction operators deterministic and
set-driven.
\end{definition}

\paragraph{Relation to Yoshinaka's finite-sample grammar.}
The fixed typing changes the present constructor in a particularly local way.
\begin{lemma}[Type guard relative to trivial typing]
\label{lem:type-guarded-subgrammar}
For every finite sample $K$, fan-out bound $f$, rank cap $r_{\max}$, and finite
typing $h$, the grammar $\mathcal B_{f,h}^{\le r_{\max}}(K)$ is obtained from
$\mathcal B_{f,h_{\mathrm{triv}}}^{\le r_{\max}}(K)$ by deleting exactly those
Rule~(Link) instances whose two tuples have different componentwise $h$-types.
The same statement holds for the uncapped operators $\mathcal B_{f,h}(K)$.
\end{lemma}
\begin{proof}
Rules~(Start), (Const), and (Comp) do not mention $h$.  Under trivial typing
every pair of same-arity tuples has equal componentwise type, so only the
equal-type test in Rule~(Link) can remove a rule.
\end{proof}

Up to notation, the trivial-typing constructor is a sample-local subgrammar of
Yoshinaka's hypothesis grammar~\cite[Section~4.1]{Yoshinaka2011}: Rule~(Start)
corresponds to his Type~III rules, Rule~(Link) to Type~II, and Rules~(Const)
and~(Comp) are instances of Type~I.  The constructors are not literally
identical.  Yoshinaka's Type~I permits any good algebraic decomposition among
observed tuple values, whereas Rule~(Comp) here requires the decomposition to
be materialized inside one observed parent occurrence.  This stricter
sample-local constructor still becomes exact after the characteristic sample
constructed below; for general $h$, Lemma~\ref{lem:type-guarded-subgrammar}
then isolates the additional finite-typing restriction.

\begin{proposition}[Fan-out-one operator equivalence]
\label{prop:fanout-one-operator}
Let $\mathcal B_h^{\mathrm{cf}}(K)$ denote the finite-sample reconstruction
operator of the companion CFG paper~\cite{kuriyamaCFG}.  For every finite
sample $K$ and every $r_{\max}\ge2$,
\[
 L\bigl(\mathcal B_{1,h}^{\le r_{\max}}(K)\bigr)
 =L\bigl(\mathcal B_{1,h}(K)\bigr)
 =L\bigl(\mathcal B_h^{\mathrm{cf}}(K)\bigr).
\]
More precisely, if $uxv\in K$ with $x\ne\eps$, then the language generated by
$[x]$ in either MCFG reconstruction equals the language generated by
$[x:u,v]$ in the CFG reconstruction.
\end{proposition}

\begin{proof}
Write $G_{\mathrm m}$ for an MCFG reconstruction and $G_{\mathrm c}$ for the
CFG reconstruction.  For the inclusion from $G_{\mathrm c}$ to
$G_{\mathrm m}$, induct on a CFG derivation.  Rule~(R1) of the companion paper
is the rank-two Rule~(Comp) obtained by splitting one observed occurrence into
two adjacent factors; Rule~(R2) only changes the recorded context and therefore
collapses to the same value state $[x]$; Rule~(R3) is Rule~(Link); and
Rules~(R4)--(R5) are simulated by Rules~(Const)--(Start).  Thus rank two already
simulates every CFG derivation.

Conversely, fix an observed occurrence $uxv\in K$ and induct on an MCFG
derivation from $[x]$.  A Rule~(Const) derivation of $x$ is reproduced by
binary Rule~(R1) splits down to letters followed by Rule~(R4).  For a
Rule~(Link) step $[x]\to[y]$, choose its witnessing common sample context
$(p,q)$.  Rule~(R2) transports $[x:u,v]$ to $[x:p,q]$, Rule~(R3) changes it to
$[y:p,q]$, and the induction hypothesis continues the derivation.  Finally,
let a Rule~(Comp) step be witnessed by a segmentation of an observed occurrence
$pxq$.  First use Rule~(R2) to move to $[x:p,q]$, then repeatedly apply
Rule~(R1) along the segmentation.  Terminal gaps are split to letters and
emitted by Rule~(R4), while every labelled child interval becomes its induced
factor--context nonterminal and is handled by the induction hypothesis.  This
simulates a composition witness of arbitrary finite rank.  The start rules,
including the direct $\eps$ rule, correspond in both constructions.  Hence the
three generated languages are equal.
\end{proof}

The CFG operator therefore retains observed occurrence contexts explicitly,
whereas the present fan-out-one operator quotients those context copies to the
factor value $[x]$; Rule~(R2) is exactly the transport that makes this quotient
language-preserving.  The value-based architecture is used here because it
extends directly to tuple values and arbitrary MCFG rule templates.

\begin{example}[Why the equal-type premise matters]
\label{ex:type-guard-purpose}
Let $L=\{\mathtt a^n\mathtt b^n\mathtt c^n:n\ge1\}$ and
$K=\{\mathtt{abc},\mathtt{aabbcc}\}\subseteq L$.  The observed factors
$(\mathtt b)$ and $(\mathtt{abbc})$ share the sample context
$\mathtt a\blank_1\mathtt c$.  With trivial typing, Rule~(Link) therefore
connects their hypothesis nonterminals.  The composition witness
$[\mathtt{aabbcc}]\to\mathtt{aa}[\mathtt b]\mathtt{bcc}$ can then use that
link to derive $\mathtt{aaabbcbcc}\notin L$.  Under $h_\Sigma^{1,1}$ the
link is blocked because $h_\Sigma^{1,1}(\mathtt b)=\mathtt b$ while
$h_\Sigma^{1,1}(\mathtt{abbc})=\langle\mathtt a,\mathtt c\rangle$.
This illustrates the local role of the equal-type premise; the strictness theorem in
Section~\ref{sec:comparison} later uses the same boundary mechanism at every
fan-out level.
\end{example}

The reconstruction is sample-consistent:
\[
 K\subseteq L\bigl(\mathcal B_{f,h}^{\le r_{\max}}(K)\bigr),
 \qquad
 K\subseteq L\bigl(\mathcal B_{f,h}(K)\bigr).
\]
Indeed, each nonempty $w\in K$ has the direct derivation
$\widehat S\to[(w)]\to(w)$ by Rules~(Start)--(Const), while
$\eps\in K$ gives the direct start rule $\widehat S\to\eps$.

These $\mathcal B$ operators are set-driven constructors, not sequential TxtEx
learners: recomputation on growing samples need not stabilize syntactically
(e.g. for successively longer $\mathtt a^n\mathtt c\mathtt b^n$).  The conservative learner below
rebuilds only when a new positive datum is rejected.

Both reconstruction operators produce finite standard MCFGs of fan-out at
most $f$ and are effectively constructible from $K,f,h$.  There are only
finitely many observed nonempty tuple occurrences.  Rule~(Link) is obtained by
grouping occurrences by concrete context and componentwise type, while
Lemma~\ref{lem:general-witness-enumeration} enumerates Rule~(Comp).  By
Lemma~\ref{lem:general-witness-enumeration}, every rule of $\mathcal B_{f,h}(K)$
has rank at most $\|K\|_+$ when $K\ne\emptyset$.

The concrete separation examples in Section~\ref{sec:comparison} also serve as
small instances of the tuple-nonterminal and composition mechanisms used here.

\subsection{Soundness}
\label{sec:soundness}

\begin{lemma}[Induced child context]
\label{lem:filling-identity}
Let $\rho$ be the template of a good rule, let
$E\in\mathcal E_e$ be a sentence context for its output tuple, and fix tuples of the arities required by the child nonterminals, with nonempty
components, for all children except child $j$.
Replacing every variable of child $j$ by its named hole and every other child
variable by the corresponding fixed component produces a context
$E_j\in\mathcal E_{d_j}$
such that
\[
 E_j[\tuple u_j]=E[\rho(\tuple u_1,\ldots,\tuple u_r)].
\]
If the sibling tuples are changed, the same construction gives the
corresponding induced context $E_j$ for that new sibling choice; the context
itself need not be identical.
\end{lemma}

\begin{proof}
Linearity and nondeletion make every variable of child $j$ occur exactly once,
and nonpermutingness puts $\xi_j^1,\ldots,\xi_j^{d_j}$ in that order.  Between two
consecutive variables of this child, nonmerging ensures either a terminal
symbol, a variable of another child, or a boundary between two output
components.  In the first case the separator is visibly nonempty; in the
second it becomes a nonempty sibling component; in the third the corresponding
internal separator of $E$ is nonempty.  Replacing the distinguished variables
by $\blank_1,\ldots,\blank_{d_j}$ therefore gives a context in
$\mathcal E_{d_j}$, and filling the holes is exactly composition with $\rho$
and then with $E$.
\end{proof}

\begin{lemma}[Composition preserves equivalence]
\label{lem:witnessed-composition}
Let $L\subseteq\Sigma^*$, let $\rho$ be the template of a good rank-$r$ rule, and let
$\tuple x_j,\tuple u_j$ have the arities required by its children and have
nonempty components.  If
$\tuple u_j\equiv_{L,h}\tuple x_j$ for every $j$, then
\[
 \rho(\tuple u_1,\ldots,\tuple u_r)
 \equiv_{L,h}
 \rho(\tuple x_1,\ldots,\tuple x_r).
\]
\end{lemma}

\begin{proof}
Homomorphic template evaluation gives equality of the two parent $h$-types,
and goodness of $\rho$ makes every parent component nonempty.  It remains to
compare distributions.  Let $E$ be an arbitrary sentence context for the
parent tuple.  If
$E[\rho(\tuple x_1,\ldots,\tuple x_r)]\in L$, replace the children one at
a time.  At step $j$, Lemma~\ref{lem:filling-identity} supplies the induced
sentence context for child $j$ under the current sibling choice.  Equality of
the child distributions transfers acceptance from $\tuple x_j$ to
$\tuple u_j$.  After all replacements,
$E[\rho(\tuple u_1,\ldots,\tuple u_r)]\in L$.  Reversing the replacements
gives the converse implication.  Since $E$ was arbitrary, the two parent
tuples have equal distributions.
\end{proof}

\begin{proposition}[Soundness of every capped hypothesis]
\label{prop:reconstruction-soundness}
Let $L$ be $(f,h)$-tuple-substitutable and $K\subseteq L$.  For every
$r_{\max}\ge1$, if $[\tuple x]\Rightarrow^*\tuple u$ in
$\widehat G_{\le r_{\max}}(K)$, then every component of $\tuple u$ is nonempty
and $\tuple u\equiv_{L,h}\tuple x$.  Consequently
\[
 L\bigl(\mathcal B_{f,h}^{\le r_{\max}}(K)\bigr)\subseteq L,
 \qquad L\bigl(\mathcal B_{f,h}(K)\bigr)\subseteq L.
\]
\end{proposition}

\begin{proof}
Induct on derivation height.  Constant rules are immediate.  For a Rule~(Link) instance $[\tuple x]\to[\tuple y]$, its defining common sample context and